\section{Въведение}
\label{sec:intro}

Уеб сървърът традиционно слуша на два порта. На порт 80 приема некриптиран HTTP, а на
порт 443 приема HTTP, защитен с TLS \cite{rfc9110}. Двата порта са две отделни входни
точки и всяка от тях знае предварително какъв протокол ще пристигне. Сървърът, който
слуша на 443, започва TLS ръкостискане веднага след приемането на връзката, без да
проверява нищо. Сървърът на 80 подава байтовете направо на анализатора на HTTP заявки.

Тази работа разглежда обратния подход. Сървърът има един-единствен слушащ сокет и за
всяка приета връзка сам решава дали тя е TLS или некриптиран HTTP. Решението се взема по
първия октет, който клиентът изпраща. Това е възможно, защото първият запис на всеки TLS
клиент започва с октет 0x16, а всяка HTTP заявка започва с главна латинска буква. Двете
множества не се пресичат, така че един октет е достатъчен.

Темата попада в дисциплината по три линии. Първо, механизмът стъпва директно на
структурата на записния слой на TLS \cite{rfc8446}: без точното значение на полето
ContentType разпознаването е невъзможно. Второ, решението „TLS или не“ се взема преди
ръкостискането, тоест в момент, когато връзката още не е защитена с нищо. Този прозорец
има собствени заплахи и те са основният предмет на анализа. Трето, след разпознаването
сървърът провежда истинско TLS ръкостискане със сертификат, верига от сертификати,
договаряне на приложния протокол (ALPN \cite{rfc7301}) и ограничения на версията.
Всяка от тези настройки има последствия за сигурността.

\subsection{Какво е направено}

Идеята произхожда от уеб сървър, който авторът разработва. За целите на проекта
механизмът е реализиран наново, самостоятелно и от нулата, в отворено хранилище с име
\code{octet} \cite{octet}. Кодът е на C++23 и използва OpenSSL. Сървърът е проверен на
Windows (MSVC, OpenSSL 3.6.4) и на Linux (GCC 14, OpenSSL 3.0.13). За macOS има
конфигурация за компилиране, но тя не е проверена и работата не твърди нищо за тази
платформа.

Конкретно проектът съдържа:
\begin{itemize}
\item сървър, който приема TLS и некриптиран HTTP/1.1 на един порт и разпознава протокола
  по първия октет;
\item пренасочване на прочетените при разпознаването октети обратно към TLS библиотеката,
  без загуба на нито един байт;
\item един общ краен срок, който покрива и разпознаването, и цялото TLS ръкостискане;
\item строго договаряне на ALPN, минимална версия TLS 1.2 или 1.3, забранено
  предоговаряне и забранена компресия;
\item собствен удостоверяващ орган за тестове и сертификат на сървъра, подписан от него;
\item автоматични тестове на ниво функция, на ниво TLS сесия в паметта и на ниво
  изпълним файл през истински сокети;
\item измерване на цената на разпознаването на две платформи.
\end{itemize}

\subsection{Какво работата твърди и какво не}

Работата твърди три неща и всяко от тях е проверено с тест или с измерване. Първо, един
октет разделя TLS от HTTP без грешка за всичките 256 възможни стойности. Второ, сървърът
с разпознаване обслужва и двата транспорта на един порт и отхвърля всичко останало.
Трето, на машината, на която е измерено, цената на разпознаването е под разделителната
способност на измерването от край до край, а в паметта тя е под една микросекунда на
връзка.

Работата не твърди, че така изграденият сървър е готов за производствена среда. В
раздел~\ref{sec:security} се показва, че самото съчетаване на защитен и незащитен
транспорт на един порт отваря път за атаки с понижаване на защитата. Това е известен
проблем, описан в препоръките на IETF \cite{rfc9325}, и работата го разглежда открито,
вместо да го заобикаля.

\subsection{Структура на работата}

Раздел~\ref{sec:basis} описва протоколната основа: записния слой на TLS, ръкостискането,
ALPN, сертификатите и началото на HTTP заявката. Раздел~\ref{sec:sharing} сравнява
известни начини за споделяне на един порт между няколко протокола. Раздел~\ref{sec:mechanism}
описва механизма в \code{octet} с откъси от кода. Раздел~\ref{sec:security} е анализът на
сигурността. Раздел~\ref{sec:verification} описва проверката, а раздел~\ref{sec:measurement}
измерването. Раздел~\ref{sec:conclusion} съдържа изводите.
